\begin{question}
    
\subsection*{Exercise 1 -- Staggered vs Un-staggered Grid}

The following code lines are taken from a larger FORTRAN code that interpolates the
temperature at a given position (described by a real pseudo-index $x_u$). The original code
works for an Arakawa C-grid. For simplicity we will consider the 1-D case where the
Arakawa grid is translated to a staggered grid. The ($x_u$) is defined using $u$-points as a
reference while the temperature is defined in an $h$-point.

\begin{verbatim}
SUBROUTINE interp(xu,T)

    ! computes temperature at the position given by xu by interpolating data
    IMPLICIT NONE
    INTEGER :: ip, im
    REAL :: xu, ax, Tint
    REAL, DIMENSION(:), INTENT(IN) :: T
    [ ... ]
    ip = NINT(xu) + 1
    im = NINT(xu)
    ax = REAL(ip) - xu
    Tint = T(im)*ax + T(ip)*(1-ax)
    [ ... ]

END SUBROUTINE
\end{verbatim}

Consider the case $x_u = 10.4$ and the temperature values:
$T(9) = 5$, $T(10) = 10$, $T(11) = 20$.

\begin{enumerate}
    \item[(a)] Using the code above compute the interpolated temperature ($T_{\text{int}}$) at $x_u$.
    
    \item[(b)] The obtained result for interpolated temperature ($T_{\text{int}}$) is wrong. Why? 
    Fix the code (explain the changes) and recalculate $T_{\text{int}}$.
    
    \item[(c)] Would the original interpolation scheme work correctly (for any $x_u$ value) if an
    unstaggered grid was used instead?
\end{enumerate}

\subsection*{Exercise 2 -- Order of Accuracy}

Consider the continuity equation of the 1-D shallow water equations:
\[
\frac{\partial h}{\partial t} = -H \frac{\partial u}{\partial x}.
\]

\begin{enumerate}
    \item[(a)] Derive the order of accuracy for the two differences of the discretised equations
    centred in time and using an unstaggered grid.
    
    \item[(b)] Repeat the previous task using instead a staggered grid.
\end{enumerate}

\end{question}
